% =============================================================================
% Chapter 1 — Introduction
% =============================================================================

\chapter{Introduction}
\label{chap:intro}

% =============================================================================
% §1.1  Context and Motivation
% =============================================================================

\section{Context and Motivation}
\label{sec:intro:context}

The global transition toward sustainable transportation has produced an
unprecedented growth in Electric Vehicle (EV) adoption. Managing the
resulting charging infrastructure presents a \emph{time-series forecasting
challenge} of considerable practical and economic magnitude: grid operators
must balance load in real time, charging networks must schedule capacity,
and energy providers must procure power days in advance. Accurate
forecasting of EV charging demand—across hourly, daily, and weekly
horizons—is therefore a critical enabling technology
\citep{acaroglu2022acn, chen2020ev}.

Classical machine learning approaches to this problem include recurrent
neural networks (RNNs), Long Short-Term Memory networks (LSTMs)
\citep{hochreiter1997lstm}, and Transformer-based models such as the
Temporal Fusion Transformer (TFT) \citep{lim2019tft}. These methods achieve
strong performance on short-term prediction but require extensive
training data, large parameter counts, and significant computational
resources for inference at scale. Echo State Networks (ESNs)
\citep{jaeger2004harnessing}—a form of reservoir computing—offer a
lightweight alternative in which only the readout layer is trained,
leaving the reservoir dynamics fixed. ESNs exploit the echo state
property to reliably encode temporal dependencies without gradient-based
learning.

Despite their appeal, ESNs operate on classical hardware, and their
representational capacity is fundamentally bounded by the classical
dynamics of their reservoir. A natural question is whether quantum
systems, whose Hilbert-space dimension grows exponentially with the
number of qubits, can provide richer temporal representations at
comparable or lower physical resource cost.

% =============================================================================
% §1.2  Quantum Reservoir Computing
% =============================================================================

\section{Quantum Reservoir Computing}
\label{sec:intro:qrc}

Quantum Reservoir Computing (QRC) \citep{jaeger2004harnessing,
fujii2017keep, appert2023quantum} proposes to replace the classical
reservoir with a quantum system: the input drives the evolution of a
quantum register, and the register's state at each timestep is measured to
produce a feature vector for a classical readout. If the quantum system
exhibits the \emph{quantum echo state property}—analogous to the echo
state property—its dynamics act as a universal black-box model for
temporal computation.

The practical appeal of QRC rests on two pillars. First, the quantum
system is \emph{not trained}; only the measurement operators and the
classical readout are optimized. This sidesteps the classical
barriers to training quantum neural networks in the NISQ era. Second,
quantum dynamics—particularly those near the edge of quantum chaos—can
generate entanglement and complex correlations that are believed to exceed
what classically simulable systems can produce efficiently
\citep{arute2019quantum, boixo2018characterizing}. This motivated a
series of experimental demonstrations showing quantum advantage for
specific temporal tasks
\citep{ Dumitrescu2018, Chen2020, Zheng2021, Ollitrault2023}.

Nevertheless, the practical advantage of QRC on real-world forecasting
tasks—particularly in the high-stakes, noisy regime of energy
infrastructure—remains an open question. Published benchmarks are
dominated by synthetic tasks (Mackey-Glass, NARMA) and small-scale
laboratory settings, leaving a significant gap between controlled
experiments and deployable forecasting systems
\citep{gosavic2024quantum, gonzalez2024benchmark}.

% =============================================================================
% §1.3  This Work
% =============================================================================

\section{This Work: QRC-EV Forecasting System}
\label{sec:intro:this_work}

This paper presents a complete \textbf{Quantum Reservoir Computing
framework for EV charging load forecasting} (QRC-EV) that bridges this
gap. Our contributions are organized along three dimensions:

\subsection{Quantum-Classical Architecture}
We implement a \textbf{stateful GPU quantum reservoir} using NVIDIA's
CUDA Quantum platform \citep{nvidiacudaq} integrated with Amazon Braket
\citep{amazonbraket}. The reservoir operates on up to 30 qubits with
input-scaled parametric gates, maintaining numerical stability across
long input sequences. The reservoir state is evolved via a
Lindblad-style dissipative channel, producing a feature map that is
subsequently processed by a classical Ridge regression readout. This
design runs entirely on cloud GPU infrastructure without the need for
physical quantum hardware access, enabling reproducible large-scale
experiments.

\subsection{Quantum Hidden Markov Model with Eligibility Traces}
Purely reservoir-based approaches treat the quantum state as a passive
temporal feature extractor. We instead pair the reservoir with a
\textbf{Quantum Hidden Markov Model (QHMM)} whose parameters are learned
via \textbf{Online Maximum Likelihood Estimation (OMLE)} subject to
complete positivity and trace preservation (CPTP) constraints. The QHMM
is represented as an \textbf{Observable Operator Model (OOM)} whose
transition structure is learned from trajectory data using a
forward-backward algorithm for state smoothing, analogous to the
Baum-Welch algorithm for classical HMMs.

Critically, we extend the QHMM-OOM framework with
\textbf{TD($\lambda$) eligibility traces} for temporal credit assignment.
The forward trace accumulates the forward-message state history;
the backward trace propagates TD errors; their product enables
efficient parameter updates without full backpropagation through time.
This allows the QHMM to learn from ongoing interaction, making it
suitable for online learning settings.

\subsection{Comprehensive Benchmarks}
We evaluate QRC-EV on four datasets spanning synthetic benchmarks
(sinusoidal, Mackey-Glass, NARMA-10) and real-world EV charging data
(ACN-Data, Palo Alto network). We compare against classical baselines
(Ridge regression, ESN, LSTM) and ablate the contribution of each
component (reservoir size, OOM smoothing, eligibility traces). Our
results demonstrate that the QHMM-OMLE pipeline achieves competitive
performance on multi-step forecasting, particularly at horizons where
the classical reservoir provides an informative prior.

% =============================================================================
% §1.4  Paper Structure
% =============================================================================

\section{Paper Structure}
\label{sec:intro:structure}

The paper is organized as follows. Chapter~\ref{chap:background}
reviews the necessary background: reservoir computing, quantum
channels and the Choi-Jamiołkowski representation, the HMM-to-OOM
connection, and the forward-backward algorithm. Chapter~\ref{chap:methods}
presents our methods: the GPU reservoir architecture, the QHMM-OOM
formalism with OMLE, the TD($\lambda$) trace mechanism, and the
optimistic planning algorithm for prediction.
Chapter~\ref{chap:experiments} describes the experimental setup,
datasets, baselines, and evaluation protocol.
Chapter~\ref{chap:results} presents and discusses the results.
Chapter~\ref{chap:conclusion} concludes with implications and open
questions.

% =============================================================================
% §1.5  Reproducibility Statement
% =============================================================================

\section*{Reproducibility Statement}
\label{sec:intro:reproducibility}

All code is publicly available at
\url{https://github.com/rezacute/medseg3d-research}. Experiments use
CUDA Quantum 0.6+ with Amazon Braket and can be run on any NVIDIA
GPU with compute capability $\geq 8.0$. Hyperparameters and random
seeds are documented in Appendix~\ref{app:hyperparameters}.
